\section{Triangular word products}%
\label{sec:asym_triangular}

The distinction between a word-level compression and a sitewise
intertwiner is already visible in elementary block matrix multiplication.

\begin{lemma}[Diagonal blocks of block-triangular products]
    \label{lem:asym_triangular}
    \lean{Matrix.blockDiag'_mul_of_blockTriangular, Matrix.blockTriangular_evalWord,
        Matrix.blockDiag'_evalWord, Matrix.blockProj_mul_blockEmbed_self,
        Matrix.blockProj_mul_blockEmbed_of_ne, Matrix.blockProj_mul_mul_blockEmbed}
    \leanok
    \uses{def:eval_word}
    Let $o$ be a finite index set with block sizes $n_k$, $k\in o$, and let
    $T^1,\ldots,T^d$ be block upper triangular for the resulting
    decomposition of $\MN{\sum_k n_k}$. Write $T^i_{kk}$ for the $k$-th
    diagonal block of $T^i$. Then $T^w$ is block upper triangular for every
    word $w$, and its $k$-th diagonal block is the word evaluation of the
    diagonal blocks,
    \begin{align}
        (T^w)_{kk} & =T^{i_1}_{kk}T^{i_2}_{kk}\cdots T^{i_N}_{kk}
        \label{eq:asym_triangular_word}
    \end{align}
    for $w=(i_1,\ldots,i_N)$. If $P_k$ is the coordinate embedding of the
    $k$-th block and $Q_k$ the coordinate projection onto it, then for every
    matrix $M$ on the same space, $Q_kMP_k=M_{kk}$, and
    \begin{align}
        Q_kP_l & =\begin{cases}\id, & k=l,\\0, & k\ne l.\end{cases}
        \label{eq:asym_block_biorthogonality}
    \end{align}
\end{lemma}

\begin{proof}
    \leanok
    For two block upper-triangular matrices $X,Y$,
    $(XY)_{kk}=\sum_l X_{kl}Y_{lk}$; the terms with $l<k$ vanish because
    $X_{kl}=0$, and those with $l>k$ vanish because $Y_{lk}=0$, leaving
    $(XY)_{kk}=X_{kk}Y_{kk}$. Induction on the word length gives
    (\ref{eq:asym_triangular_word}) and the persistence of block
    triangularity. The identity $Q_kMP_k=M_{kk}$ and
    (\ref{eq:asym_block_biorthogonality}) are direct computations from the
    coordinate description of $P_k$ and $Q_k=P_k^{\mathsf T}$.
\end{proof}

\begin{remark}[Periodic equality does not imply a sitewise intertwiner]
    \label{rem:asym_zipper}
    Take the one-dimensional normal tensor $A^1=(1)$, $A^2=(2)$, and the
    bond-dimension-three tensor
    \begin{align}
        B^1
         & =\begin{pmatrix}0&1&0\\0&1&0\\0&0&0\end{pmatrix},
         &
        B^2
         & =\begin{pmatrix}0&0&0\\0&2&1\\0&0&0\end{pmatrix}.
        \notag
    \end{align}
    Both matrices are upper triangular with diagonal blocks $0$, $A^i$, $0$.
    With $P=e_2$ and $Q=e_2^{\mathsf T}$, Lemma~\ref{lem:asym_triangular}
    gives $QB^wP=A^w$ and $QP=1$ for every nonempty word $w$, so
    $\tr(B^w)=A^w=\tr(A^w)$: the periodic vectors of $A$ and $B$ agree at
    every positive length.

    There is nevertheless no nonzero right sitewise intertwiner: a vector
    $v=(x,y,z)^{\mathsf T}$ with $B^1v=v$ and $B^2v=2v$ would need $x=y$,
    $z=0$ from the first equation and $x=0$ from the first coordinate of the
    second, forcing $v=0$. Symmetrically, no nonzero left intertwiner
    $uB^1=u$, $uB^2=2u$ exists: the first equation gives $u_1=u_3=0$ while
    the third coordinate of the second gives $u_2=2u_3=0$. Periodic equality
    can therefore force a multiplicative subquotient while admitting neither
    one-sided sitewise intertwiner, so it cannot imply a biorthogonal zipper
    pair. This is why the theorem below is stated as a compression
    $W_sB^wV_s=C_s^w$ rather than as an eigenvector relation.
\end{remark}

\section{The word algebra and normal tensors as simple modules}%
\label{sec:asym_word_algebra}

\begin{definition}[Word algebra and positive ideal]
    \label{def:asym_word_algebra}
    \lean{WordAlgebra, WordAlgebra.aug}
    \leanok
    Fix $d\geq1$. The \emph{word algebra} is the free unital associative
    algebra $F=\C\langle x_1,\ldots,x_d\rangle$, realized as the algebra of
    finite $\C$-linear combinations of words in the $d$ letters under
    concatenation, so that the words form a basis. The \emph{augmentation}
    $\mathrm{aug}:F\to\C$ is the algebra map sending every generator to
    zero; on a basis word $w$, $\mathrm{aug}(w)=1$ if $w$ is empty and $0$
    otherwise. The \emph{positive ideal} is the two-sided ideal
    $F_+=\ker(\mathrm{aug})$ of elements with zero constant term.
\end{definition}

\begin{definition}[Word-algebra module of a tensor]
    \label{def:asym_word_module}
    \lean{MPSTensor.WordModule}
    \leanok
    \uses{def:asym_word_algebra, def:mps_tensor, def:eval_word}
    A tensor $C=(C^i)_{i=0}^{d-1}$ on $\C^D$ defines an $F$-module structure
    on $\C^D$ through the representation $\rho_C(x_i)=C^i$, extended
    multiplicatively so that a word $w$ acts by $\rho_C(w)=C^w$.
\end{definition}

\begin{definition}[Word traces]
    \label{def:asym_word_traces}
    \lean{WordAlgebra.traceChar, WordAlgebra.traceWord}
    \leanok
    \uses{def:asym_word_algebra}
    For a finite-dimensional $F$-module $M$, the \emph{trace character} is
    the linear functional $p\mapsto\tr(\rho_M(p))$ on $F$, and the
    \emph{word trace} of a word $w$ is $\tr(\rho_M(w))$.
\end{definition}

\begin{lemma}[Word traces of a tensor]
    \label{lem:asym_word_traces_tensor}
    \lean{MPSTensor.traceWord_wordModule}
    \leanok
    \uses{def:asym_word_traces, def:asym_word_module}
    On the word-algebra module of a tensor $C$, the word trace of $w$ is
    the periodic coefficient $\tr(C^w)$.
\end{lemma}

\begin{proof}
    \leanok
    \uses{def:asym_word_module}
    The action of $w$ is conjugate to the matrix $C^w$ under the
    identification of the module with $\C^D$, and the trace of a linear
    map is the trace of its matrix.
\end{proof}

\begin{lemma}[Normal tensors define simple modules]
    \label{lem:asym_normal_simple}
    \lean{MPSTensor.isSimpleModule_wordModule_of_isNormal}
    \leanok
    \uses{def:asym_word_module, def:normal}
    If $C$ is normal on $\C^D$ with $D>0$, its word-algebra module is
    simple.
\end{lemma}

\begin{proof}
    \leanok
    Let $K\subseteq\C^D$ be a nonzero invariant subspace and $u\in K$
    nonzero. Normality gives a length $L$ for which the span of the
    length-$L$ word evaluations is all of $\MN{D}$; since $K$ is invariant
    under every such word and hence under their span, $Xu\in K$ for every
    $X\in\MN{D}$. Every vector is $Xu$ for a suitable $X$ because $u\neq0$,
    so $K=\C^D$.
\end{proof}

Multiplying every matrix of a tensor by a nonzero scalar does not change its
invariant subspaces, so a weighted block $\mu_kA_{(k)}$ built from a normal
tensor $A_{(k)}$ and $\mu_k\in\C^\times$ remains a simple $F$-module.

\section{Separating simple modules}%
\label{sec:asym_separation}

\begin{lemma}[Positive-word character lemma]
    \label{lem:asym_character}
    \notready
    \uses{def:asym_word_module}
    Let $U,V$ be finite-dimensional $F$-modules with
    $\tr(\rho_U(p))=\tr(\rho_V(p))$ for every $p\in F_+$. Then the
    semisimplifications of $U$ and $V$ contain every simple $F$-module with
    the same multiplicity, except possibly the one-dimensional module
    $\mathbf0$ on which every generator acts as zero: there are integers
    $a,b\geq0$ and a semisimple module $M$ without a $\mathbf0$ summand such
    that $U^{\mathrm{ss}}\cong M\oplus\mathbf0^{\oplus a}$ and
    $V^{\mathrm{ss}}\cong M\oplus\mathbf0^{\oplus b}$.
    % The source proof decomposes the finite-dimensional image algebra
    % modulo its Jacobson radical by Artin--Wedderburn and compares
    % traces of central idempotents.
\end{lemma}

\begin{lemma}[Positive-ideal separation of linear functionals]
    \label{lem:asym_ker_aug_separation}
    \lean{WordAlgebra.linearMap_eq_on_ker_aug}
    \leanok
    \uses{def:asym_word_algebra}
    Two $\C$-linear functionals on $F$ agreeing on every nonempty word
    agree on the positive ideal $F_+$.
\end{lemma}

\begin{proof}
    \leanok
    Their difference agrees with $(\chi_1(1)-\chi_2(1))\cdot\mathrm{aug}$ on
    every basis word, by induction on the monoid-algebra presentation, hence
    on all of $F$ by linearity; on $F_+$ the augmentation vanishes, so the
    two functionals agree there.
\end{proof}

\begin{lemma}[The zero module is the only simple module with trivial
        generator action]
    \label{lem:asym_zero_module}
    \lean{WordAlgebra.finrank_eq_one_of_isSimpleModule_of_forall_smul_eq_zero}
    \leanok
    \uses{def:asym_word_algebra}
    A finite-dimensional simple $F$-module on which every generator acts as
    zero is one-dimensional. It is the module $\mathbf0$ of
    Lemma~\ref{lem:asym_character}.
\end{lemma}

\begin{proof}
    \leanok
    If every generator acts as zero, $F$ acts through the augmentation, so
    the $F$-submodules and the $\C$-subspaces coincide; simplicity as a
    $\C$-vector space forces dimension one.
\end{proof}

\begin{theorem}[Burnside's theorem at module level]
    \label{thm:asym_burnside}
    \lean{WordAlgebra.exists_smul_eq_of_isSimpleModule}
    \leanok
    \uses{def:asym_word_module}
    Every $\C$-linear endomorphism of a finite-dimensional simple
    $F$-module $Q$ is the action of some element of $F$.
\end{theorem}

\begin{proof}
    \leanok
    By Schur's lemma over the algebraically closed field $\C$, the
    commutant $\End_F(Q)$ consists of scalars, so $Q$ is finite-dimensional
    over its commutant. The Jacobson density theorem then makes the action
    map $F\to\End_{\C}(Q)$ surjective.
\end{proof}

\begin{theorem}[Separating a simple module from a finite non-isomorphic
        family]
    \label{thm:asym_separation}
    \lean{WordAlgebra.exists_aug_eq_zero_smul_eq_self,
        WordAlgebra.exists_aug_eq_zero_smul_eq_self_smul_eq_zero,
        MPSTensor.exists_aug_eq_zero_of_forall_isEmpty_linearEquiv}
    \leanok
    \uses{def:asym_word_module, def:asym_word_algebra}
    Let $Q$ be a finite-dimensional simple $F$-module on which some
    generator acts nonzero. There is $u\in F_+$ acting as the identity on
    $Q$. If $T$ is a finite-dimensional simple module not isomorphic to
    $Q$, there is $p\in F_+$ acting as the identity on $Q$ and as zero on
    $T$. More generally, if $C_s$, $s\in S$, is a finite family of
    word-algebra modules of tensors, each simple and none isomorphic to
    $Q$, there is $p\in F_+$ acting as the identity on $Q$ and as zero on
    every $C_s$.
\end{theorem}

\begin{proof}
    \leanok
    \uses{thm:asym_burnside}
    For the first claim, Burnside's theorem realizes, for a basis
    $\{b_j\}$ of $Q$ and a nonzero vector $x_i\cdot q_0$, the maps sending
    $x_i\cdot q_0\mapsto b_j$ and $q\mapsto(\text{coordinate }j\text{ of
        }q)\cdot q_0$; assembling them into
    $u=\sum_j\alpha_jx_i\beta_j$ gives an element of $F_+$ acting as the
    identity on $Q$. For the pairwise separation, Schur's lemma shows every
    $F$-linear endomorphism of $Q\times T$ preserves the two
    non-isomorphic simple factors, so $Q\times T$ is finite-dimensional
    over its (scalar) commutant; Jacobson density on $Q\times T$ produces
    $a\in F$ acting as the identity on the first factor and zero on the
    second, and $p=au$ acts as the identity on $Q$ and as zero on $T$. The
    finite-family statement follows by induction on $S$, multiplying the
    separators for the successive members.
\end{proof}

\begin{lemma}[Trace additivity over an invariant submodule]
    \label{lem:asym_trace_additivity}
    \lean{WordAlgebra.traceWord_eq_add_quotient}
    \leanok
    \uses{def:asym_word_traces}
    For a finite-dimensional $F$-module $M$ and a submodule $K\leq M$, the
    word trace of $M$ is the sum of the word traces of $K$ and of $M/K$.
\end{lemma}

\begin{proof}
    \leanok
    This is the additivity of the ordinary trace of an endomorphism over an
    invariant subspace, applied to the action of each word on $M$.
\end{proof}

\begin{theorem}[Nilpotency from vanishing trace powers]
    \label{thm:asym_nilpotent_trace_powers}
    \lean{LinearMap.isNilpotent_of_forall_trace_pow_eq_zero}
    \leanok
    An endomorphism of a finite-dimensional complex vector space, all of
    whose positive powers have vanishing trace, is nilpotent.
\end{theorem}

\begin{proof}
    \leanok
    The characteristic polynomial is determined by the power sums of its
    roots (Newton's identities), so vanishing trace powers force it to
    equal $X^{\dim}$; the Cayley--Hamilton theorem then gives nilpotency.
\end{proof}

\begin{theorem}[Identification of a simple quotient with a target block]
    \label{thm:asym_identification}
    \lean{MPSTensor.exists_linearEquiv_wordModule_of_isSimpleModule_quotient}
    \leanok
    \uses{def:asym_word_traces, def:asym_word_module}
    Let $M$ be a finite-dimensional $F$-module whose word traces are, on
    every nonempty word, the sum of the word traces of a finite family of
    simple tensor modules $C_s$, $s\in S$. Let $K\leq M$ have simple
    quotient $M/K$ on which some generator acts nontrivially. Then $M/K$ is
    isomorphic to $C_s$ for some $s\in S$.
\end{theorem}

\begin{proof}
    \leanok
    \uses{thm:asym_separation, lem:asym_ker_aug_separation,
        thm:asym_nilpotent_trace_powers}
    If $M/K$ were isomorphic to no $C_s$, Theorem~\ref{thm:asym_separation}
    supplies $p\in F_+$ acting as the identity on $M/K$ and as zero on
    every $C_s$. Since the word trace of $M$ agrees with the sum of the
    word traces of the $C_s$ on every nonempty word, and $p\in F_+$,
    Lemma~\ref{lem:asym_ker_aug_separation} shows that every positive power
    of the action of $p$ on $M$ has vanishing trace, so this action is
    nilpotent by Theorem~\ref{thm:asym_nilpotent_trace_powers}. A nilpotent
    action of $p$ on $M$ descends to a nilpotent action on the quotient
    $M/K$, contradicting that $p$ acts there as the identity on a nonzero
    module.
\end{proof}

\section{The composition-series flag}%
\label{sec:asym_flag}

\begin{lemma}[Composition-series triangular gauge]
    \label{lem:asym_composition_gauge}
    \lean{exists_linearEquiv_blockTriangular_of_flag}
    \leanok
    \uses{def:asym_word_module}
    Let $B=(B^i)_i$ act on $\C^{D_B}$, and choose a composition series
    $0=H_0\subset H_1\subset\cdots\subset H_r=\C^{D_B}$ of its
    word-algebra module. There is $G\in\GL_{D_B}(\C)$ such that every
    $T^i=GB^iG^{-1}$ is block upper triangular for the resulting
    decomposition into $r$ blocks, and the $k$-th diagonal block of $T^i$
    is the action induced by $B^i$ on the simple subquotient
    $H_k/H_{k-1}$.
\end{lemma}

\begin{proof}
    \leanok
    Choose a basis of $H_1$ and extend it successively to bases of
    $H_2,\ldots,H_r$. In the resulting basis, every $H_k$ is a coordinate
    initial segment; since every $H_k$ is invariant under every $B^i$, all
    matrices $B^i$ are block upper triangular in this basis, with $k$-th
    diagonal block the induced action on $H_k/H_{k-1}$.
\end{proof}

\begin{definition}[Labelled invariant flag]
    \label{def:asym_flag_data}
    \lean{MPSTensor.FlagData}
    \leanok
    \uses{def:asym_word_module}
    Let $M$ be a finite-dimensional $F$-module, $S$ a finite set of target
    slots, and $C_s$, $s\in S$, a family of tensors. A \emph{labelled
        invariant flag} of $M$ over $(S,C)$ is an invariant chain
    $0=H_0\leq H_1\leq\cdots\leq H_r=M$ together with a bijective labelling
    of the steps $\{1,\ldots,r\}$ by $S\sqcup\{1,\ldots,z\}$, for some
    $z\geq0$, such that a step labelled by $s\in S$ has subquotient
    $H_k/H_{k-1}$ isomorphic to the word-algebra module of $C_s$, and a
    step labelled by an index in $\{1,\ldots,z\}$ has one-dimensional
    subquotient on which every generator acts as zero.
\end{definition}

\begin{lemma}[A family of normal blocks with zero total word trace is
        empty]
    \label{lem:asym_flag_base_case}
    \lean{MPSTensor.finset_eq_empty_of_forall_sum_trace_evalWord_eq_zero}
    \leanok
    \uses{def:normal, def:asym_word_traces}
    If the word traces of a finite family of normal tensors of positive
    bond dimension sum to zero on every nonempty word, the family is
    empty.
\end{lemma}

\begin{proof}
    \leanok
    \uses{thm:asym_separation, thm:asym_nilpotent_trace_powers}
    Otherwise Theorem~\ref{thm:asym_separation} applied to one member $C_{s_0}$
    of the family supplies $u\in F_+$ acting as the identity on $C_{s_0}$;
    vanishing total word trace forces every positive power of the action of
    $u$ on the product of the modules in the family to have trace zero, so
    this action is nilpotent by Theorem~\ref{thm:asym_nilpotent_trace_powers}.
    A nilpotent action of $u$ contradicts its action as the identity on the
    nonzero summand $C_{s_0}$.
\end{proof}

\begin{theorem}[The flag theorem]
    \label{thm:asym_flag_theorem}
    \lean{MPSTensor.exists_flagData}
    \leanok
    \uses{def:asym_flag_data, def:normal, def:asym_word_traces}
    Let $C_s$, $s\in S$, be a finite family of normal tensors of positive
    bond dimension, and let $M$ be a finite-dimensional $F$-module whose
    word traces agree, on every nonempty word, with the sum of the word
    traces of the $C_s$. Then $M$ carries a labelled invariant flag over
    $(S,C)$.
\end{theorem}

\begin{proof}
    \leanok
    \uses{lem:asym_flag_base_case, lem:asym_zero_module,
        thm:asym_identification, lem:asym_trace_additivity}
    Induction on $\dim M$. If $M=0$, Lemma~\ref{lem:asym_flag_base_case}
    forces $S=\emptyset$, and the empty chain is the required flag. If
    $M\neq0$, choose a maximal proper invariant submodule $K\leq M$ (a
    coatom of the invariant-submodule lattice, which exists because $M$ is
    finite-dimensional), so $M/K$ is simple. If some generator acts
    nontrivially on $M/K$, Theorem~\ref{thm:asym_identification} identifies
    $M/K$ with some block $C_{s_0}$; Lemma~\ref{lem:asym_trace_additivity}
    shows the word trace of $K$ agrees with the sum over the remaining
    blocks $S\setminus\{s_0\}$, and the induction hypothesis (applied to
    $K$, of strictly smaller dimension) supplies a flag of $K$, extended by
    the matched final step $K\subset M$ labelled $s_0$. If instead every
    generator acts as zero on $M/K$, Lemma~\ref{lem:asym_zero_module} shows
    $M/K$ is one-dimensional; Lemma~\ref{lem:asym_trace_additivity} shows
    the word trace of $K$ already agrees with the sum over all of $S$, and
    the induction hypothesis supplies a flag of $K$, extended by the
    unmatched final step $K\subset M$.
\end{proof}

\section{The multi-block asymmetric compression theorem}%
\label{sec:asym_main}

\begin{theorem}[Multi-block asymmetric compression theorem]
    \label{thm:asym_multiblock}
    \lean{MPSTensor.MultiBlockCompression,
        MPSTensor.exists_multiBlockCompression_of_isNormal,
        MPSTensor.MultiBlockCompression.dim_eq,
        MPSTensor.MultiBlockCompression.isReduction,
        MPSTensor.MultiBlockCompression.left_mul_right_of_ne,
        MPSTensor.MultiBlockCompression.evalWord_remainder_eq_zero}
    \leanok
    \uses{def:same_mpv2_pos, def:normal}
    Let $A=\bigoplus_{s\in S}C_s$ be a finite direct sum of nonzero normal
    tensors $C_s=(C_s^i)_i$ of bond dimension $D_s$, with repetitions
    allowed, and let $B=(B^i)_i$ be an arbitrary tensor of bond dimension
    $D_B$. Assume that the periodic vectors of $A$ and $B$ agree at every
    positive length,
    $\ket{V^{(N)}(B)}=\ket{V^{(N)}(A)}$ for every $N>0$; equivalently,
    $\tr(B^w)=\tr(A^w)$ for every nonempty word $w$. Then there are an
    integer $z\geq0$, an integer $r=|S|+z$, an invertible matrix
    $G\in\GL_{D_B}(\C)$, and maps $V_s:\C^{D_s}\to\C^{D_B}$,
    $W_s:\C^{D_B}\to\C^{D_s}$, with the following properties.
    % Plain enumerate with explicit tags: the kernel key-value list syntax
    % [label=(\roman*)] requires LaTeX2e >= 2025, which CI does not have.
    \begin{enumerate}
        \item[(i)] Every $GB^iG^{-1}$ is block upper triangular for one common
              decomposition into $r$ blocks.
        \item[(ii)] Exactly $|S|$ of the diagonal slots are matched with the
              target slots: writing $j(s)$ for the position matched with $s$,
              there is $X_s\in\GL_{D_s}(\C)$ with
              $\bigl(GB^iG^{-1}\bigr)_{j(s),j(s)}=X_sC_s^iX_s^{-1}$.
        \item[(iii)] The remaining $z$ diagonal slots are one-dimensional zero
              tensors.
        \item[(iv)] The compression pairs are mutually biorthogonal: $W_sV_t=
                  \id_{D_s}$ if $s=t$ and $0$ otherwise.
        \item[(v)] For every nonempty word $w$, $W_sB^wV_s=C_s^w$.
        \item[(vi)] Writing $\widehat C^i=\sum_{s\in S}V_sC_s^iW_s$ and
              $R^i=B^i-\widehat C^i$, every product
              $R^{i_1}R^{i_2}\cdots R^{i_r}=0$: the non-unital algebra generated
              by the $R^i$ is nilpotent.
        \item[(vii)] $D_B=\sum_{s\in S}D_s+z$.
    \end{enumerate}
\end{theorem}

\begin{center}
    % Formula: clause (v), W_sB^wV_s=C_s^w for every nonempty word w.
    % Source: clauses (iv)-(v) of Theorem~\ref{thm:asym_multiblock}.
    % Ink-to-index map: the two end rings are the compression maps $W_s$
    % (left) and $V_s$ (right); the interior nodes are copies of $B$,
    % with upper legs carrying the letters of $w$.
    % Contraction set: every horizontal virtual bond between consecutive
    % copies of $B$ and between each $B$ and its adjacent ring is
    % contracted; the physical legs and the two outer virtual legs (the
    % $D_s$-dimensional indices of the matrix $C_s^w$) remain open.
    % Boundary signature: (virtual-west=1 | virtual-east=1 |
    % physical-up=2 | physical-down=0) for the finite schematic shown,
    % with the two explicit lettered nodes carrying the drawn physical
    % ports.
    \tenkzkernel
    \begin{tenkz}[rows={wire}, cols=5, west=open, east=open]
        \tn[skin=ring]{W_s} &
        \tn[label pos=s, ports={90:physical:$i_1$}]{B} &
        \tn[skin=dots]{} &
        \tn[label pos=s, ports={90:physical:$i_N$}]{B} &
        \tn[skin=ring]{V_s}
    \end{tenkz}
    \quad $=$ \quad
    \begin{tenkz}[rows={wire}, cols=3, west=open, east=open]
        \tn[label pos=s, ports={90:physical:$i_1$}]{C_s} &
        \tn[skin=dots]{} &
        \tn[label pos=s, ports={90:physical:$i_N$}]{C_s}
    \end{tenkz}
    \\[2ex]
    % Formula: clause (iv), W_sV_t=\id_{D_s} if s=t and 0 otherwise.
    % Source: clause (iv) of Theorem~\ref{thm:asym_multiblock}.
    % Ink-to-index map: the two rings are the compression maps $W_s$ and
    % $V_t$, joined with no intervening tensor; the two outer virtual
    % legs are the $D_s$- and $D_t$-dimensional indices of the matrix
    % $W_sV_t$.
    % Contraction set: the single virtual bond between the two rings.
    % Boundary signature: (virtual-west=1 | virtual-east=1 |
    % physical-up=0 | physical-down=0).
    \begin{tenkz}[rows={wire}, cols=2, west=open, east=open]
        \tn[skin=ring]{W_s} & \tn[skin=ring]{V_t}
    \end{tenkz}
    \quad $=\;\delta_{st}\,\id$
\end{center}

\begin{proof}
    \leanok
    \uses{thm:asym_flag_theorem, lem:asym_composition_gauge,
        lem:asym_triangular}
    By Theorem~\ref{thm:asym_flag_theorem}, applied with $M$ the
    word-algebra module of $B$, whose word traces agree with the sum of
    those of the $C_s$ by hypothesis, $\C^{D_B}$ carries a labelled
    invariant flag over $(S,C)$: an invariant chain $0=H_0\leq\cdots\leq
        H_r=\C^{D_B}$ whose subquotients are, in some order, the $C_s$ (each
    once) and $z=r-|S|$ copies of the zero module. This chain is in
    particular a composition series, so Lemma~\ref{lem:asym_composition_gauge}
    supplies $G\in\GL_{D_B}(\C)$ for which $T^i=GB^iG^{-1}$ is block upper
    triangular with these subquotients as diagonal blocks, proving (i)--(iii).

    Let $P_{j(s)}$, $Q_{j(s)}$ be the coordinate embedding and projection
    of the matched slot and set $V_s=G^{-1}P_{j(s)}X_s$,
    $W_s=X_s^{-1}Q_{j(s)}G$. Then $W_sV_t=X_s^{-1}Q_{j(s)}P_{j(t)}X_t$, and
    (\ref{eq:asym_block_biorthogonality}) gives (iv). For a nonempty word
    $w$, $B^w=G^{-1}T^wG$, and Lemma~\ref{lem:asym_triangular} gives
    $Q_{j(s)}T^wP_{j(s)}=X_sC_s^wX_s^{-1}$, so
    $W_sB^wV_s=X_s^{-1}Q_{j(s)}T^wP_{j(s)}X_s=C_s^w$, proving (v).

    For the remainder, $G\widehat C^iG^{-1}=\sum_{s\in S}P_{j(s)}X_sC_s^iX_s^{-1}Q_{j(s)}$
    is the block-diagonal matrix agreeing with $T^i$ on every matched
    diagonal block; every unmatched diagonal block of $T^i$ is already
    zero, so $GR^iG^{-1}=T^i-G\widehat C^iG^{-1}$ is strictly block upper
    triangular for the $r$-block decomposition. A nonzero entry of a
    product of $r$ strictly block upper-triangular $r$-block matrices would
    need a strictly increasing chain of $r+1$ block indices, which does not
    exist among $r$ positions; hence $GR^{i_1}\cdots R^{i_r}G^{-1}=0$,
    proving (vi). Finally, dimensions add along the composition series, and
    each zero-module step contributes dimension one, giving (vii).
\end{proof}

\begin{theorem}[Weighted canonical-form compression]
    \label{thm:asym_weighted_compression}
    \lean{MPSTensor.exists_weightedMultiBlockCompression_of_isNormal,
        MPSTensor.MultiBlockCompression.left_mul_evalWord_mul_right_smul}
    \leanok
    \uses{def:same_mpv2_pos, def:normal}
    Let $A=\bigoplus_{k=1}^g\bigoplus_{q=1}^{n_k}\mu_kA_{(k)}$ be a
    canonical target with each $A_{(k)}$ normal of bond dimension $D_k$
    and $\mu_k\in\C^\times$, and suppose
    $\ket{V^{(N)}(B)}=\ket{V^{(N)}(A)}$ for every $N>0$. For every block
    label $k$ and copy $q\in\{1,\ldots,n_k\}$, there are maps
    $V_{k,q}:\C^{D_k}\to\C^{D_B}$ and $W_{k,q}:\C^{D_B}\to\C^{D_k}$,
    mutually biorthogonal, with
    \begin{align}
        W_{k,q}B^{i_1}\cdots B^{i_N}V_{k,q}
         & =\mu_k^NA_{(k)}^{i_1}\cdots A_{(k)}^{i_N}
        \label{eq:asym_weighted_compression}
    \end{align}
    for every $N\geq1$ and every word.
\end{theorem}

\begin{proof}
    \leanok
    \uses{thm:asym_multiblock}
    Apply Theorem~\ref{thm:asym_multiblock} to the slots
    $C_{k,q}=\mu_kA_{(k)}$: for a word of length $N$,
    $C_{k,q}^{i_1}\cdots C_{k,q}^{i_N}=\mu_k^NA_{(k)}^{i_1}\cdots
        A_{(k)}^{i_N}$, so clause (v) of Theorem~\ref{thm:asym_multiblock}
    gives (\ref{eq:asym_weighted_compression}).
\end{proof}

% The Lean formalization of Theorem~\ref{thm:asym_multiblock} takes the
% hypothesis in the word-trace form $\tr(B^w)=\tr(A^w)$, and its conclusion
% chooses every matching isomorphism $X_s$ in clause (ii) to be the identity,
% which is stronger than the printed clause above; the underlying
% construction already produces this normalization because the composition
% series identifies each matched subquotient with the block's own
% word-algebra module, not merely with an isomorphic copy.

\begin{proposition}[Periodic data of a compression]
    \label{prop:asym_compression_trace}
    \lean{MPSTensor.MultiBlockCompression.trace_evalWord_eq_sum,
        MPSTensor.MultiBlockCompression.mpv_eq_sum}
    \leanok
    \uses{thm:asym_multiblock, lem:asym_triangular}
    Conversely, whenever a tensor $B$ admits the block triangular form of
    Theorem~\ref{thm:asym_multiblock}, clauses (i)--(iii), for a family
    $(C_s)_{s\in S}$, its periodic data are those of the direct sum:
    $\tr(B^w)=\sum_{s\in S}\tr(C_s^w)$ for every nonempty word $w$, that
    is, $\ket{V^{(N)}(B)}=\ket{V^{(N)}(\bigoplus_sC_s)}$ for every $N\geq1$.
\end{proposition}

\begin{proof}
    \leanok
    \uses{lem:asym_triangular}
    The trace is invariant under the change of basis and is the sum of the
    traces of the diagonal blocks; by Lemma~\ref{lem:asym_triangular} the
    diagonal blocks of $B^w$ are the words of the diagonal blocks, which are
    the $C_s^w$ on the matched slots and $0$ on the one-dimensional zero
    slots.
\end{proof}

\begin{proposition}[A vanishing remainder gives sitewise intertwiners]
    \label{prop:asym_split_intertwiners}
    \lean{MPSTensor.MultiBlockCompression.mul_right_eq_right_mul,
        MPSTensor.MultiBlockCompression.left_mul_eq_mul_left}
    \leanok
    \uses{thm:asym_multiblock}
    In the setting of Theorem~\ref{thm:asym_multiblock}, suppose that the
    remainder vanishes, $R^i=0$ for every physical index $i$. Then for
    every $i$ and every $s\in S$,
    \begin{align}
        B^iV_s & =V_sC_s^i, &
        W_sB^i & =C_s^iW_s.
    \end{align}
\end{proposition}

\begin{proof}
    \leanok
    Since $R^i=0$, $B^i=\sum_{t\in S}V_tC_t^iW_t$. Multiplying on the
    right by $V_s$, or on the left by $W_s$, and using the
    biorthogonality $W_tV_s=\delta_{ts}\id_{D_s}$ of clause (iv) leaves
    only the term $t=s$.
\end{proof}

\begin{proposition}[Projector-weighted direct sums]
    \label{prop:asym_projector_weighted_sum}
    \lean{MPSTensor.MultiBlockCompression.ofProjectorSum,
        MPSTensor.MultiBlockCompression.remainder_ofProjectorSum}
    \leanok
    \uses{thm:asym_multiblock, lem:asym_triangular}
    Let the bond space of $B$ be $\C^m\otimes\C^{D}$ and let $\Pi_1,\ldots$
    be mutually orthogonal rank-one idempotents of $\C^m$ cut out by an
    invertible matrix $U$ (the idempotents $U^{-1}E_{jj}U$). If every
    letter has the form
    $B^i=\sum_{s\in S}\mu_s\,\Pi_{j(s)}\otimes C_s^i$ with distinct
    positions $j(s)$, then $B$ admits the block triangular form of
    Theorem~\ref{thm:asym_multiblock} onto the weighted blocks $\mu_sC_s$,
    with $z=(m-|S|)D$ zero factors, the remainder $R^i=0$, and sitewise
    intertwiners $B^iV_s=V_s\,\mu_sC_s^i$ and $W_sB^i=\mu_sC_s^iW_s$ for
    every slot: the extension splits.
\end{proposition}

\begin{proof}
    \leanok
    \uses{lem:asym_triangular, prop:asym_split_intertwiners}
    Conjugating by $U\otimes\id$ turns each $\Pi_{j(s)}$ into a matrix
    unit, so every letter becomes block diagonal with the blocks
    $\mu_sC_s^i$ at the positions $j(s)$ and zero blocks elsewhere; block
    diagonality gives the triangular form and the vanishing remainder, and
    Proposition~\ref{prop:asym_split_intertwiners} gives the intertwining
    relations.
\end{proof}

